\documentclass[12pt]{article}

% Title information
\title{Understanding Tax Enforcement: Questionnaire (and Experimental Design)}
\author{Tim Kircali and Matthias Weber}
\date{\today} % Use \date{} for no date

\begin{document}

\maketitle
%Abstract
\tableofcontents
\section{Treatments}
\subsection{Inequality Priming}
Is K.TS really necessary? If people don't know anything at all about the tax system and tax brackets and so on, 

\section{Questionnaire}
The structure of the questionnaire is shown in the table of contents.
\subsection{Title Page / Consent Form}
\subsection{Background Questions}
These questions have not been formulated year, but I think we should cover the following background information. 
\begin{enumerate}
    \item Gender?
    \item Age?
    \item Total household income before taxes?
    \item Born in (US/CH/..)?
    \item zip code?
    \item marital status?
    \item children?
    \item ethnicity/race?
    \item level of education?
    \item employment status?
    \item social welfare program? food stamp? cash transfer?
    \item social class? (lower class poor, working classs, middle class, upper-middle class, upper class))
    \item where on liberal conservative spectrum?
    \item political affiliation?
    \item did you vote, if yes for who?
    \item do you follow the news?
    \item are you informed about economic policy?
    \item Why is it important to be informed?
    \item which sources do you use to stay informed?
\end{enumerate}
\subsection{Information Treatments}
Experimental treatments. These are not defined yet, first I am trying to focus on the questionnaire. 
\subsection{Knowledge}
Overall, these questions measure knowledge and perceptions of relevant issues. What we need to ask ourselves is whether we are unintentionally priming people with issues in a way that influences their political reasoning. If we prime people with issues, we may need to experimentally vary whether they are primed. 
\subsubsection{Knowledge about inequality}
This block focuses on general knowledge of inequality. 
\begin{enumerate}
    \item How has the share of total U.S. income that goes to the top 1\% in the U.S. evolved over the past 30 years
    \begin{itemize}
        \item It has increased by a lot; 
        \item It has increased somewhat; 
        \item It has remained the same; 
        \item It has decreased somewhat; 
        \item It has decreased by a lot
    \end{itemize} $\rightarrow$ This question asks about knowledge/perception of how inequality has changed?
    \item $\rightarrow$ We could ask a similar question about wealth inequality.
    \item $\rightarrow$ We could ask a similar question about intergenerational mobility
\end{enumerate}
%(We could consider asking them about inequality in connection with taxes. A question which arises is, however, if we prime them about inequality by asking them. So would they really reason their policy views with inequality, if we had not talked about that?)
\subsubsection{Knowledge about the tax system in general}
This block concentrates on general knowledge regarding tax policy. The questions are very Similar as in Stantchevas "Understanding Tax Policy" Paper. 
\begin{enumerate}
    \item Imagine a middle class household that is right at the middle of the income distribution, such that half of all households in the U.S. earn more than this household and half earn less. What share of their income do you think such a household pays in (federal) income taxes? \\
    $\rightarrow$ This question is used to elicit the belief about the median payers tax rate.
    \item What is the (fedderal) top tax rate (in percent)? \\
    $\rightarrow$ To elicit the belief/knowledge of the top tax rate. 
    \item For how many people does the top tax rate apply? \\
    $\rightarrow$ To elicit the belief/knowledge of the amount of people in top tax bracket. 
    \item Can you now try to describe the complete tax schedule by adjusting the interactive tool. (While you're doing that, you have to keep the government revenue constant.) \\
    $\rightarrow$ To elicit the belief/knowledge of the complete tax schedule. However, I am not sure if we should give them the information about the government revenue.  
    \item For the following questions assume there is no tax evasion (and tax avoidance), and that everyone pays their taxes as the law requires them to. \\
    $\rightarrow$ This part is added to disentangle tax evasion effects on the effective tax schedule. By doing so, we are giving people an indication that tax evasion affects progressivity, so maybe we have to find another way to disentangle between tax schedule progressivity in general and how it is connected to tax evasion. Maybe by experimentally vary the schedule. And I am not sure yet if we should exclude tax avoidance too. 
    \begin{enumerate}
        \item Do you think that, broadly speaking, everyone in the U.S. (in your country) currently pays approximately the same share of income in federal personal income taxes, or do you think people pay very different shares of income in (federal) personal income taxes?" (Stantcheva, 2021)
        \item (If response to (a) not “Everyone pays more or less...”) Do you think that people with higher incomes pay a higher or lower share of their total income in (federal) personal income taxes than people with lower incomes?" (Stantcheva, 2021)
    \end{enumerate} $\rightarrow$ With these questions we elicit their beliefs/understanding of the progression in the tax schedule. 
\end{enumerate}
\subsubsection{Knowledge about evading inequality}
This section concentrates on knowledge about tax evasion and tax evasion inequalitites. 
\begin{enumerate}
    \item For the following questions, you don't have to assume that tax evasion (and avoidance) doesn't exist. Think about the current situation of tax evasion and how people actually pay tax.   \\
    $\rightarrow$ After asking the question from above, we now have to introduce again that tax evasion and avoidance exists. By doing so, we are giving people an indication that tax evasion affects progressivity, so maybe we have to find another way to disentangle between tax schedule progressivity in general and how it is connected to tax evasion. Maybe by experimentally vary the schedule. 
    \begin{enumerate}
        \item Do you think that, broadly speaking, everyone in the U.S. (in your country) currently pays approximately the same share of income in federal personal income taxes, or do you think people pay very different shares of income in (federal) personal income taxes?" (Stantcheva, 2021)
        \item (If response to (a) not “Everyone pays more or less...”) Do you think that people with higher incomes pay a higher or lower share of their total income in (federal) personal income taxes than people with lower incomes?" (Stantcheva, 2021)
    \end{enumerate} $\rightarrow$ With these questions we elicit their beliefs/understanding of the progression in the tax system, including effects of tax evasion and avoidance. 
    \item Now assume Person A has higher income than person B. Both can deduct 10000.- from their taxable income. Which of the two persons will save more in paid taxes? \\ $\rightarrow$ This question tests how well people understand the concept of marginal tax rates in tax enforcement, i.e. whether they understand that a reduction in taxable income translates into higher savings for people in higher tax brackets.
    \item What do you think, how much of tax revenue is getting lost due to tax evasion (in percent)? \\$\rightarrow$ With this question we analyse their perception about the amount of income that is evaded. However, I don't know yet, what kind of measures we will use to evaluate if they are right.
    \item What do you think, broadly speaking, would you say that the level of tax evasion is related to the level of income? What do you think, do people that earn more, also evade more or less? \\$\rightarrow$ With this question we could analyse their perception about the distribution of tax evasion. Generally, we could also ask them if it is related to other stuff.  Note: Is this maybe rather a mechanism question?
    \item What do you think, which of the following groups evades the most? 
    \begin{itemize}
        \item Large businesses
        \item Small businesses
        \item High-income households
        \item Middle-income households
        \item Low-income households
    \end{itemize}$\rightarrow$ With this question we could analyse their perception about who is evading. Note: Is this maybe rather a mechanism question?
\end{enumerate}
\subsubsection{Knowledge about enforcement policies}
This section concerns knowledge about enforcement policies; About audit probabilities and also punishments.
\begin{enumerate}
    \item Overall, what do you think, how many (percent of) tax declarations are being closely examined/audited? \\$\rightarrow$ This questions tests for general knowledge/perception of enforcement level.
    \item What are the audit probabilities of (the following) subgroups? 
    \begin{itemize}
        \item Large businesses
        \item Small businesses
        \item High-income households
        \item Middle-income households
        \item Low-income households
    \end{itemize} $\rightarrow$ This question tests if people think certain subgroups are being treated differently.
    \item Do you think that chances of being audited are related to income?  Do you think audit probabilities are higher/lower for higher incomes? \\ $\rightarrow$ These questions test if people think there is a link between higher enforcement and income. 
    \item Do you think that chances of being audited are related to other (/one of the following) things? \\ $\rightarrow$ This question test if people see any connection between the level of enforcement and other variables. The variables I have not defined yet. 
    \item What do you think, how do people get punished when evading? \\ $\rightarrow$ This question tries to understand what people's knowledge/perception about punishments of tax enforcement is. However, the question is really unprecise yet. One idea would be to measure their beliefs related to other punishments. As example, if they think tax evasion of 10000 francs is punished more heavily than crime X. 
\end{enumerate}
\subsubsection{Knowledge about automatic exchange of information policies}
This section asks about knowledge of AEI policies. 
\begin{enumerate}
    \item Do you know what automatic exchange of information is? \\ $\rightarrow$ 
    With this question we would first ask if they even know about the existence of AEI. But probably the question is not well formulated yet, because people might know about it but don't know the word. 
    \item What would you say to which of the following and how much personal financial information have tax authorities access to? \\$\rightarrow$ With this question I am trying to understand their current knowledge/beliefs of the level of AEI. However, I have not defined yet, from which information options they can choose. 
    \item Do you know which kind of information (which of the following information) is transferred in AEI? \\$\rightarrow$ This question asks about what kind of information are actually transmitted in AEI, but we would have to define more precisely of which AEI policy we are talking about. 
    \item Have you heard about pre-filed tax declarations? \\ $\rightarrow$ I am asking this to understand if people know about pre-filed tax declarations, however, maybe the question is not really good, because people might know about the existence, but not about the name. 
    \item Did you know that in other countries taxes are automatically deducted from salaries? \\ $\rightarrow$ Maybe there is a better way to ask this question, but the aim is to understand if they are aware that in other tax systems the timing of when taxes are paid and by whom is different.
    \item Did you know that pre-filed tax declarations significantly reduce effort to file tax declarations? \\$\rightarrow$ Obviously, this question is not really good. But the goal is to test if people know that this reduces effort. Maybe this is however also a mechanism question. 
\end{enumerate}

\subsection{Views}
The goal of this section is to elicit people's general views (views about government, fairness views, views about others), that might be somehow influence their reasoning about policy views. But generally, I might not have not enough view questions. So this is a part I could still very much improve.  
\subsubsection{Government Views}
The following questions have been copied almost 1:1 from Stantcheva (Understanding Trade).
\begin{enumerate}
    \item How much of the time do you think you can trust our (federal) government to do what is right? \\$\rightarrow$ This question measures trust in government. 
    \item Some people think the government is trying to do too many things that should be left to individuals and businesses. Others think that government should do more to solve our country’s problems. Which come closer to your own view?\\$\rightarrow$ This question measures the preferred level of state intervention. 
    \item Where would you rate yourself on a scale of 1 to 5, where 1 means you think the government should do only those things necessary to provide the most basic government functions, and 5 means you think the government should take active steps in every area it can to try and improve the lives of its citizens?\\$\rightarrow$ This question measures the preferred level of state intervention. 
    \item Of every tax dollar that goes to the federal government (in Washington, D.C.), how many cents would you say are wasted?\\$\rightarrow$ This measures the view how wasteful the government is.
    \item Are you very satisfied, somewhat satisfied, somewhat dissatisfied, or very dissatisfied with the way the (federal) government (in Washington) is dealing with the problems the country is facing today? \\ $\rightarrow$ This question measures if people think that government is doing a good job or not. 
    \item Consider now a list of functions the federal government could serve.  On a 1 to 5 scale, please say how much responsibility you think the government should have for each — with 1 meaning the government should have no responsibility at all and 5 meaning the government should have total responsibility in this area:
    \begin{itemize}
        \item Reducing income differences between the rich and the poor.
        \item Reducing the transmission of wealth from one generation to the other.
        \item Making sure Americans have adequate health care.
        \item Reducing the differences in opportunities between children from wealthy and poor families.
        \item Regulating trade to and from the U.S. to protect American producers and consumers.
        \item Maintaining a stable financial system and ensuring that credit markets work.
        \item Ensuring a stable dollar.
        \item Providing a minimum standard of living for all.
    \end{itemize} $\rightarrow$ With these questions we are eliciting the views on what governments should do.  
\end{enumerate}
\subsubsection{Equity, Inequality, Fairness Views}
Probably I could work a bit more on these questions below...I need some more relevant view questions and maybe also structure them a bit better. 
\begin{enumerate}
    \item Generally, would you say that the current tax system is unfair? 
    \item How should the optimal tax brackets look like? Can you now try to describe the complete tax schedule by adjusting the interactive tool. (While you're doing that, you have to keep the government revenue constant.) \\$\rightarrow$ This question should elicit the preferred level of progression/redistribution within the tax system. 
    \item Generally, would you say people can be trusted to hand in their tax declaration honestly? \\ $\rightarrow$ With this question we are trying to elicit if people think that other people can be trusted to hand in their tax declaration honestly.
    \item Let's say there is an easy way for a taxpayer to evade tax, is it the taxpayer's responsibility to file honestly or is it the government's fault that there is such a way?  \\ $\rightarrow$ With a question like this I trying to understand if people think it is okay to evade if there is the option.  
    \item Let's say there is an easy way for a taxpayer to evade tax, is it okay if a taxpayer makes use of this option and evades?  \\ $\rightarrow$ With a question like this I trying to understand if people think it is okay to evade if there is the option.
    \item Which has more to do with why a person is rich?  Because she or he worked harder than others; Because she or he had more advantages than others $\rightarrow$ With this question I am trying to understand equality of opportunity. 
\end{enumerate}

\subsection{Mechanism Questions}
The following questions help us to understand how people reason about policy views. We ask them about their perceived effects/mechanisms policies bring. Mechanism questions could also be asked with a experimental variation like in Stantcheva (2021), with a "neutral","I/Household similar to mine","Women formulation" . \\ $\rightarrow$ Sometimes it is a little bit difficult to say if these questions are rather a "knowledge", a "view", or a "mechanism" question. 

\subsubsection{On inequality within the tax system}
With the following questions we are trying to understand what people think how tax enforcement is connected and affected with inequality, especially if certain subgroups have higher chances to evade and also if generally the tax system is a tool to tackle inequality.
\begin{enumerate}
    \item What do you think, which of these subgroups have the most possibilities/options to evade and avoid?
     \begin{itemize}
        \item Large businesses
        \item Small businesses
        \item High-income households
        \item Middle-income households
        \item Low-income households
    \end{itemize} $\rightarrow$ With this question we want to understand what people think which of the taxpayer types have the most options to evade 
    \item And which of these subgroups would you say evades the most? 
    \begin{itemize}
        \item Large businesses
        \item Small businesses
        \item High-income households
        \item Middle-income households
        \item Low-income households
    \end{itemize} $\rightarrow$ With this question we try to elicit their perception about which of the subgroups actually evades the most.
    \item Who would you say avoids the most? 
    \begin{itemize}
        \item Large businesses
        \item Small businesses
        \item High-income households
        \item Middle-income households
        \item Low-income households
    \end{itemize} $\rightarrow$ With this question we try to elicit their perception about which of the subgroups actually avoids the most. (Again, here we have the problem that we don't know if these persons can disentangle between avoidance and evasion)
    \item Do you think the design of tax system and tax enforcement is an important tool to tackle income inequality? \\$\rightarrow$ This questions asks if the tax system is an important tool to tacke inequality. Generally, it would be nice to disentangle between the design of tax brackets \& tax law and the design of tax enforcement.  
    \item Assume that the level of tax evasion increases; Would you say income and wealth inequality is affected by that? If yes, in which direction? \\ $\rightarrow$ We are trying to understand if people think that higher tax enforcement has an impact on inequality. 
\end{enumerate}
\subsubsection{On the effects of increased tax enforcement}
These questions elicit the perceived effects of increased tax enforcement. 
\begin{enumerate}
    \item Do you think you would gain/suffer financially from increased tax enforcement (=more audits)? \\ $\rightarrow$ With this question we try to understand if people think that they, personally, would be affected by changes in tax enforcement. This question could be asked in several experimental variations; "neutral", "I/Household similar to mine","Women formulation" 
    \item Do you think, overall, tax evasion is reduced if we invest more in tax enforcement (=more audits)? \\ $\rightarrow$ With this question we try to understand the effectiveness of tax enforcement on evasion. 
    \item What do you think, which of these subgroups evades the most taxes?
    \begin{itemize}
        \item Large businesses
        \item Small businesses
        \item High-income households
        \item Middle-income households
        \item Low-income households
    \end{itemize} $\rightarrow$ With this question we try to understand what they think which of the subgroups evade the most. (We have the same question in the knowledge section, so I am not sure where it belongs)
    \item Who do you think would gain/suffer the most from increased tax enforcement (=more audits)?
    \begin{itemize}
        \item Large businesses
        \item Small businesses
        \item High-income households
        \item Middle-income households
        \item Low-income households
    \end{itemize} $\rightarrow$ To elicit what people think whou would be the major loosres of a enforcement policy. 
    \item Generally, do you think tax revenue would increase in that sense that it would (over) compensate investments in tax enforcement if we invest more in tax enforcement (=more audits)? \\ $\rightarrow$ This question elicits the belief about the efficiency of tax enforcement policy. 
    \item Do you think increased tax enforcement would reduce income inequality? \\ $\rightarrow$ Elicits perceived effects of tax enforcement on income inequality. (We have the same question above, not sure where it belongs)
    \item Do you think increased tax enforcement would reduce wealth inequality? \\ $\rightarrow$ Elicits perceived effects of tax enforcement on wealth inequality. 
    \item Do you think an increase in tax enforcement makes the tax system more fair overall? \\ $\rightarrow$ Elicits perceived effects of tax enforcement on fairness of tax system. 
    \item Which punishments would you say are the most effective (regardless of whether they are appropriate) to prevent tax evasion? Evaluate the effectiveness of the following punishments on reducing tax evasion. \\ $\rightarrow$ Generally, this question would ask people to evaluate the effectiveness of a selection of punishments.
    \item Which punishments would you say are the most appropriate (regardless of whether they are effective) to prevent tax evasion? Evaluate the appropriateness of the following punishments on reducing tax evasion. \\ $\rightarrow$ In general, this question would ask people how appropriate certain punishments would be. Note: I am trying here to disentangle how appropriate a punishment is from how effective it is. However, it may be difficult to disentangle effectiveness from appropriateness, as appropriateness could also be justified by effectiveness. 
    \item Consider taxpayers/evaders in income class XY, which of the following taxpayers punishments would be the most effective to reduce tax evasion? \\ $\rightarrow$ This question helps us to understand if people think different punishments work better for different subgroup of evaders. 
    \item Consider taxpayers/evaders in income class XY, which of the following taxpayers punishments would be the most appropriate to reduce tax evasion? \\ $\rightarrow$ This question helps us to understand if people think different punishments are more appropriate for different subgroup of evaders. 
\end{enumerate}

\subsubsection{On the effects of increased automatic exchange of information}
These questions elicit the perceived effects of increased AEI. 
\begin{enumerate}
    \item Do you think you would gain/suffer financially (=have to pay more taxes) from increased automatic exchange of information with tax authorities? \\ $\rightarrow$ This question elicits the perceived financial gain/loss of automatic exchange of information. It also could be asked in different Variations: "neutral", "I/Household similar to mine","Women formulation" 
    \item Generally, do you think tax evasion is reduced if we increase automatic exchange of information? \\ $\rightarrow$ This question elicits the perceived effects on tax evasion if we increase AEI. 
    \item Which group would you expect to pay more taxes if we increased automatic information exchange?
    \begin{itemize}
        \item Large businesses
        \item Small businesses
        \item High-income households
        \item Middle-income households
        \item Low-income households
    \end{itemize} $\rightarrow$ This question elicits the perceived effects on tax evasion of certain subgroups if we increase AEI.
    \item Do you think tax revenue would increase if we increase automatic exchange of information? \\ $\rightarrow$ This question asks for the perceived effects of AEI on tax revenue.
    \item Do you think tax revenue would would increase such that it would compensate investments of aei? \\ $\rightarrow$ This question asks if the introduction of AEI would compensate costs. 
    \item Do you think income inequality would increase/decrease if we increased automatic exchange of information? \\ $\rightarrow$ Tests for perceived effect of AEI on income inequality. 
    \item Do you think wealth inequality would increase/decrease if we increased automatic exchange of information? \\ $\rightarrow$ Tests for perceived effect of AEI on wealth inequality. 
    \item Do you think automatic exchange of information makes the tax system more fair overall? \\ $\rightarrow$ Tests for perceived effect of AEI on fairness of tax system.
    \item Generally, do you think that increasing automatic exchange would bring opportunities, such that as example  tax payments could be automatically deducted from wage payments? \\ $\rightarrow$ Asks for perceived opportunities of AEI.
    \item Generally, do you think that increasing automatic exchange would bring opportunities, such that as example tax filing costs could be reduced? (General, You variation, Someone else Variation) \\ $\rightarrow$ Asks for perceived opportunities of AEI.
    \item Which group would you expect would profit from reduced tax filing costs if we increased automatic information exchange?
    \begin{itemize}
        \item Large businesses
        \item Small businesses
        \item High-income households
        \item Middle-income households
        \item Low-income households
    \end{itemize} $\rightarrow$ This question asks more detailed for perceived opportunities of AEI for certain subgroups.
    \item Generally, do you think that increasing automatic exchange would bring opportunities, such that as example tax control costs could be reduced? \\ $\rightarrow$ This question asks for perceived opportunities of AEI for tax administrators.
    \item For which of the following groups would you expect tax administrations to benefit from reduced tax control costs if we increased automatic information exchange?
    \begin{itemize}
        \item Large businesses
        \item Small businesses
        \item High-income households
        \item Middle-income households
        \item Low-income households
    \end{itemize} $\rightarrow$ This question aims to understand in which subgroups AEI could help tax administartors.
    \item Generally, do you think it is a problem if the Government (and tax authorities) collect data of individuals and firms? \\ $\rightarrow$ With this question we try to understand if people generally think it is a problem if governments collect data.
    \item Evaluate for the following groups how okay it would be if the Government (and tax authorities) would collect data?
        \begin{itemize}
        \item Large businesses
        \item Small businesses
        \item High-income households
        \item Middle-income households
        \item Low-income households
    \end{itemize} $\rightarrow$ With this question we try to understand if people think it is okay for governments to collect data for certain subgroups.
    \item Do you think people/firms would loose data privacy if more information is transferred automatically to tax authorities? (General, You variation, Someone else Variation) \\ $\rightarrow$ With this question we try to understand if people think they (or other people or firms) would loose data privacy with the introduction of AEI, compared to before
    \item Do people mind about tax authorities collecting their data? (General, You variation, Someone else Variation) \\ $\rightarrow$ With this question we want to elicit if they actually mind (or if they think other people/firms mind) when the Government collects their data. 
\end{enumerate} 

\subsection{Policy Views}
Finally, these questions ask about concrete policy views. 
\subsubsection{on tax enforcement}
\begin{enumerate}
    \item Generally, would you want that enforcement activities / audits are increased (which would also mean higher enforcement costs)? \\ $\rightarrow$ This measures general policy view for increased enforcement.
    \item Which taxpayers should audits focus on? 
     \begin{itemize}
        \item Large businesses
        \item Small businesses
        \item High-income households
        \item Middle-income households
        \item Low-income households
    \end{itemize} $\rightarrow$ This questions measures where people want tax administrators to focus.
    \item What is your preferred tax audit probability schedule? How would you distribute a budget of tax audits among these subgroups? 
    \begin{itemize}
        \item Large businesses
        \item Small businesses
        \item High-income households
        \item Middle-income households
        \item Low-income households
    \end{itemize} $\rightarrow$ This question does the same as above, but more precise.
    \item Should the major goal of tax audits be to maximize tax revenue or should it rather be to increase fairness and uncover fraud? \\ $\rightarrow$ With this question I am trying to elicit what the preferred goal is. Note: I am not sure if I should just differentiate between the maximizing tax revenue goal and other goals, or if I also should differentiate between other goals. 
    \item Should tax audits main goal be to maximize tax revenue and therefore focus on tax evaders that weight the heaviest or should tax audits mainly reduce tax evasion overall, increasing fairness also for small evaders? \\ $\rightarrow$ This question is very similar to the one above, however it explicitly asks about on which taxpayers it should focus. 
    \item Generally, should tax audits be targeted or should they be random? \\ $\rightarrow$ elicits if people support targeted tax audits or if they perceiva that as unfair.  
    \item Generally, which fines would you support?  $\rightarrow$ (Here, I still have to work a bit on it)
 \end{enumerate}
\subsubsection{on automatic exchange of information \& Privacy}
\begin{enumerate}
    \item Generally, would you support that automatic exchange of information of tax data would be introduced? \\ $\rightarrow$ Elicits general support of AEI. 
    \item How would you evaluate your support for automatic exchange of information for the following subgroups?
        \begin{itemize}
        \item Large businesses
        \item Small businesses
        \item High-income households
        \item Middle-income households
        \item Low-income households
    \end{itemize} $\rightarrow$ differentiates for support of AEI in subgroups.
    \item For which kind of data would you support automatic exchange of information: 
    \begin{itemize}
        \item Wage data.
        \item Health insurance data.
        \item Bank data.
        \item Payment data. (For the VAT)
    \end{itemize} $\rightarrow$ differentiates for support of AEI in different types of data.
    \item Would you support the state to offer an optional automatic exchange of information?  \\ $\rightarrow$ I think this question makes not much sense, because we cannot really disentangle between something. But generally it would be helpful for government if there is a optional option.
    \item If there is an optional AEI, would you use it? 
    \item Would you support tax payments being automatically deducted from your wage? \\ $\rightarrow$ Maybe this is not really a policy View question.
    \item Would you support tax authorities offering pre-filled tax declarations?\\ $\rightarrow$ Maybe this is not really a policy View question, but it elicits if people want pre-filled tax declarations.
    \item How much would you be willing to pay for pre-filled tax declarations?\\ $\rightarrow$ This question would measure how much they were willing to pay for it. 
\end{enumerate}

\subsection{(Willingness to pay for information)}
\subsection{(Self-reported questions about attention during study)}
\subsection{(Feedback about the survey)}
\end{document}
